\documentclass[11pt]{article}
\usepackage{amsmath,amssymb,amsthm}
\usepackage[margin=1in]{geometry}
\usepackage[hidelinks]{hyperref}
\newtheorem{lemma}{Lemma}
\newtheorem{theorem}{Theorem}
\begin{document}
\setlength{\emergencystretch}{1em}
\title{Sector separation for support $\{4,6\}$ at every positive Gamma shape}
\author{}\date{}\maketitle
For $\beta>0$ let $\mathcal L_\beta(X^j)=(\beta)_j$.  Put
\[
 T_{n,\gamma}(t)=\sum_{j=0}^n\binom nj(\gamma)_{2j}t^j,\qquad
 P_n(t)=T_{n,4n+\beta}(t).
\]
We prove that $P_n$ and $P_{n+1}$ have no common complex root for
every positive integer $n$, without an arithmetic condition on $\beta$.

The proof combines an exact shape-shift-four adjoint with a strict
Pick quotient and a Gamma-product zero-free sector.  Each imported
statement is used below within its exact domain.

In a fixed pair set $\gamma=4n+\beta$, $M=\gamma+2n=6n+\beta$,
$s=2(n+1)$, and
\[
 B=B_{n,\gamma},\quad C=B_{n+1,\gamma},\quad
 B=B_0(u)+zB_1(u),\quad C=C_0(u)+zC_1(u),\quad u=z^2.
\]
Here the auxiliary recurrence is
\[
 B_{0,\gamma}=1,\qquad
 B_{j+1,\gamma}=(z+2j+\gamma+1)B_{j,\gamma}-zB_{j,\gamma}'.
\]
Write
\[
 A=(M+2)(M+3)(M+4)(M+5),\quad
 k=(M+4)(M+5),\quad b=s(2M+7),
\]
\[
 U(u)=A+bu,\qquad W(u)=su(k+u).
\]
The exact shape-shift condition used below is that every common
root $t$ gives $u=-1/t\ne0$ with
\[
 U(u)C_i(u)+W(u)B_i(u)=0,\qquad i=0,1.                 \tag{1}
\]

\begin{lemma}[A necessary angle]\label{lem:four-sector-angle}
For $n\ge2$, a nonreal common root $t$ satisfies $\Re t>0$ and
\[
 \cos^2(\arg t)>Q_n(M),\qquad
 Q_n(M)=\frac{2(n+1)(2M+7)}{(M+2)(M+3)}.              \tag{2}
\]
The function $Q_n(M)$ is strictly decreasing on $M>0$.
\end{lemma}
\begin{proof}
Since $\gamma>0$ and $\gamma\ge n-2$, the
\href{LOWER_FOUR_QUADRATIC_GAP_SFC.tex}
{general shift-four Lemma \texttt{lem:offset-four-transform}}
gives (1).  Here $\gamma=4n+\beta>2n-2$, so
\href{LOWER_FIVE_QUADRATIC_SECTOR_SFC.tex}
{Lemma \texttt{lem:pick-general}} applies.
In particular $B_1$ has only negative real zeros, and
\[
 \frac{\Im(C_1/B_1)(u)}{\Im u}>0
\]
for every nonreal $u$.  The nonreal point $u=-1/t$ is not a zero
of $B_1$ or $U$: the only zero of $U$ is the negative real point
$-A/b$.  Thus (1) gives $C_1/B_1=-W/U$ there.

For $u=x+iy$, an exact expansion gives
\[
 \frac{\Im(W(u)\overline{U(u)})}{y}
       =s\{Ak+2Ax+b(x^2+y^2)\}.
\]
The expression in braces is strictly negative at a common root.
Substituting $u=-1/t$ and multiplying by $|t|^2$ yields
\[
 Ak|t|^2-2A\Re t+b<0.                                \tag{3}
\]
All three constants are positive.  Hence $\Re t>0$.
Regard (3) as a quadratic inequality in $|t|$, with
$\Re t=|t|\cos(\arg t)$.  Its discriminant must be positive,
so $\cos^2(\arg t)>bk/A=Q_n(M)$.

Differentiation gives
\[
 Q_n'(M)=
 -\frac{2(n+1)(2M^2+14M+23)}{(M+2)^2(M+3)^2}<0.
\]
\end{proof}

\begin{lemma}[Four exact sector comparisons]\label{lem:four-sector-bounds}
Let $n\ge2$ and $\beta>0$.  Every point satisfying (2) lies in a
zero-free sector of $T_{n,4n+\beta}$.
\end{lemma}
\begin{proof}
Put
\[
 \mathcal R_n(a)=\frac{4(n-1)(a+1)}{(a-n+1)^2},
 \qquad a=\gamma/2.
\]
It is strictly decreasing for $a>n-1$.  By
Lemma~\ref{lem:four-sector-angle}, $Q_n$ is strictly decreasing in $M$.
The Gamma-product sector
\href{LOWER_FIVE_QUADRATIC_SECTOR_SFC.tex}
{Lemma \texttt{lem:moment-sector}} says that, for any
$q>\mathcal R_n(a)$, the lower polynomial
has no zero in
\[
 |\arg t|<\Omega(q),\qquad
 \Omega(q)=\pi-2\arctan\sqrt q.
\]
All applications have $\gamma>2n-2$.  The following exact margins
give the four useful choices:
\[
\begin{aligned}
 8-\mathcal R_n(2n)
   &=\frac{4(5n+3)}{(n+1)^2}>0,\\
 \frac{144}{25}-\mathcal R_n(9n/4)
   &=\frac{16(485n+244)}{25(5n+4)^2}>0,\\
 3-\mathcal R_n(3n)
   &=\frac{20n+7}{(2n+1)^2}>0,\\
 1-\mathcal R_n(6n)
   &=\frac{n^2+30n+5}{(5n+1)^2}>0.
\end{aligned}                                                     \tag{4}
\]
Let $N=2(n+1)(2M+7)$ and $D=(M+2)(M+3)$.  Direct expansion gives
\[
\begin{aligned}
 \bigl[N-\tfrac8{13}D\bigr]_{M=13n/2}&=\frac{2(130n+67)}{13}>0,\\
 \bigl[N-\tfrac12D\bigr]_{M=8n}&=26n+11>0,\\
 \bigl[N-\tfrac27D\bigr]_{M=14n}&=\frac{2(175n+43)}7>0.
\end{aligned}                                                     \tag{5}
\]
Use the decreasing functions $Q_n$ and $\mathcal R_n$ as follows.

If $0<\beta\le n/2$, then $M\le13n/2$ and $a\ge2n$.
Consequently $Q_n(M)>8/13$ and $q=8$ is permitted.
Since
\[
 \cos\Omega(8)=\frac79,\qquad
 \frac8{13}-\left(\frac79\right)^2=\frac{11}{1053}>0,
\]
(2), together with $\Re t>0$, forces $|\arg t|<\Omega(8)$.

If $n/2\le\beta\le2n$, then $M\le8n$, $a\ge9n/4$,
$Q_n(M)>1/2$, and $q=144/25$ is permitted.  Here
\[
 \cos\Omega(144/25)=\frac{119}{169},\qquad
 \frac12-\left(\frac{119}{169}\right)^2
                    =\frac{239}{57122}>0.
\]
Again the necessary angle lies in the zero-free sector.

If $2n\le\beta\le8n$, then $M\le14n$, $a\ge3n$,
$Q_n(M)>2/7>1/4$, and $q=3$ gives
$\Omega(3)=\pi/3$.  The same implication follows.

Finally, if $\beta\ge8n$, then $a\ge6n$ and $q=1$ gives
the zero-free open right half-plane.  This contains every point
allowed by Lemma~\ref{lem:four-sector-angle}.
The four ranges cover all positive real $\beta$.
\end{proof}

\begin{theorem}\label{thm:gamma-lower-four-quadratic}
For every positive real $\beta$ and every integer $n\ge1$,
$T_{n,4n+\beta}$ and $T_{n+1,4(n+1)+\beta}$ are coprime over
$\mathbb C$.  Every nonzero complex phase $X^4(A_0+B_0X^2)$
satisfies Gamma-$\beta$ Strong Factorial at every consecutive
positive pair.
\end{theorem}
\begin{proof}
There is no real common root: the even-index member has the
strictly positive integral
\[
 T_{j,4j+\beta}(t)=\frac1{\Gamma(4j+\beta)}
    \int_0^\infty x^{4j+\beta-1}(1+tx^2)^j e^{-x}\,dx
       >0\qquad(j\ {\rm even},\ t\in\mathbb R).
\]
For $n=1$, the lower polynomial is linear with a real root,
so this suffices.  For $n\ge2$, a nonreal common root would satisfy
Lemma~\ref{lem:four-sector-angle}, contradicting
Lemma~\ref{lem:four-sector-bounds}.

If $A_0B_0\ne0$, the exact moment normalization is
\[
 \mathcal L_\beta\bigl((X^4(A_0+B_0X^2))^j\bigr)
 = A_0^j(\beta)_{4j}T_{j,4j+\beta}(B_0/A_0).
\]
The omitted factors are nonzero.  If one coefficient vanishes,
the phase is a nonzero monomial, whose positive moments are nonzero.
\end{proof}

\paragraph{Scope and evidence.}
This proof covers all positive real shapes, including
transcendental shapes.  It asserts no unrestricted binomial SFC
theorem or half-plane stability of the lower moment polynomial.
The written sector argument supplies every index and shape
quantifier.  Symbolic controls check the identities, angular margins
and normalizations; finite checks are not used as an index theorem.
Internal mathematical audit does not constitute independent
whole-proof replay, formal verification or external review.
\end{document}
